% реферат по теме 'Технологии создания и обработки аудио'
\documentclass{SibFU-docs}

% доп. пакеты
\usepackage{graphicx}
\usepackage{float}
\usepackage{hyperref}
\usepackage{caption}
\usepackage{subcaption}
\usepackage{amsmath}
% предотвартим ошибку с \Bbbk перед загрузкой пакета amssymb
\let\Bbbk\relax
\usepackage{amssymb}
\usepackage{booktabs}
\usepackage{tabularx}
\usepackage{longtable}
\usepackage{enumitem}
\usepackage{multirow}

% гиперссылки
\hypersetup{
    colorlinks=true,
    linkcolor=black,
    citecolor=blue,
    urlcolor=blue
}

% библиография
\addbibresource{report.bib}

% титульный лист
\begin{document}
\setlist[itemize]{left=1.3cm, itemsep=1pt, topsep=2pt} % для выравнивания списков
\setlist[enumerate]{left=1.3cm, itemsep=1pt, topsep=2pt}
\renewcommand{\contentsname}{}
\mycovertitle
{Гуманитарный институт}
{Кафедра прикладной информатики в искусстве и интерактивном медиа}
{Технологии создания и обработки аудио}
{ст. преподаватель И.~Р.~Нигматуллин}
{Логинова Ю.В, Захаров И.М., Мельников С.В., Тужикова В.В.,\\Тетерина П.А., Тугаринова А.И., Черепанова А.А.}
 
\begin{center}
\bfseries РЕФЕРАТ
\end{center}
\noindent

Реферат по теме: «Технологии создания и обработки аудио» содержит 29 страниц и 20 использованных источников.

Ключевые слова: АУДИОСИГНАЛ, ДИСКРЕТИЗАЦИЯ, КОДЕК, СЖАТИЕ ДАННЫХ, ПСИХОАКУСТИКА, АУДИОМОНТАЖ, ПРОСТРАНСТВЕННЫЙ ЗВУК, АУДИОФОРМАТ, ПЕРЕДАЧА АУДИОДАННЫХ, ИММЕРСИВНЫЙ ЗВУК

В работе исследуются фундаментальные технологии создания, обработки и представления цифрового аудио. Рассмотрены базовые принципы преобразования аналогового звукового сигнала в цифровую форму, включая процессы дискретизации и квантования, лежащие в основе любого аудиотехнологического стека. Особое внимание уделено методам кодирования, сжатия и передачи аудиоданных, анализу различий между сжатием без потерь и с потерями, а также ключевым форматам файлов и областям их применения. Проанализирован современный инструментарий для редактирования и монтажа звука в цифровой среде, а также перспективные направления развития технологий объемного и пространственного звучания, определяющие будущее аудиоиндустрии.

Основные выводы:

\begin{enumerate}
    \item Цифровое представление звука, базирующееся на теореме Котельникова-Найквиста (дискретизация) и процессе квантования, является технологическим фундаментом для записи, обработки и распространения аудиоконтента, обеспечивая его точность, воспроизводимость и возможность цифровой манипуляции.
    
    \item Эффективность хранения и передачи аудио напрямую зависит от применяемых алгоритмов сжатия. Сжатие без потерь (например, FLAC) сохраняет полную идентичность оригиналу и критически важно в профессиональной сфере, тогда как сжатие с потерями (MP3, AAC), основанное на психоакустических моделях маскирования, обеспечивает радикальное уменьшение размера файла за счет незаметного для слуха упрощения данных, что доминирует в потребительском сегменте.

    \item Создание и тонкая обработка аудиоконтента осуществляется в специализированных цифровых аудио-рабочих станциях (DAW), которые предоставляют комплекс инструментов для записи, редактирования, сведения и мастеринга, превращая исходные записи в законченные музыкальные или звуковые произведения.

    \item Будущее аудиотехнологий связано с переходом к иммерсивным и интерактивным форматам. Технологии объемного и пространственного звука, такие как Dolby Atmos и Ambisonics, создают трёхмерную звуковую сцену, что кардинально меняет опыт восприятия в кинематографе, музыке, видеоиграх и виртуальной реальности, открывая новые возможности для творчества и коммуникации.
\end{enumerate}

Рекомендации:

\begin{itemize}[label=-]
    \item Для долгосрочного архивирования и профессиональной работы с исходными материалами использовать форматы без потерь (WAV, FLAC). Для распространения финального контента конечному потребителю, особенно через онлайн-платформы и потоковые сервисы, применять современные эффективные форматы сжатия с потерями (AAC, Opus).

    \item При освоении навыков аудиомонтажа и обработки начинать с изучения базовых операций в DAW (обрезка, склейка, нормализация, применение простых эффектов), постепенно переходя к более сложным процессам сведения и мастеринга для достижения коммерческого качества звука.

    \item Учитывать растущий тренд на пространственный звук при создании нового аудиоконтента для кино, игр и иммерсивных медиа. Изучать основы работы с соответствующими технологиями и форматами для сохранения актуальности создаваемых продуктов.

    \item При выборе кодеков и форматов для проектов всегда оценивать целевую платформу воспроизведения и каналы распространения, обеспечивая максимальную совместимость и оптимальное соотношение между качеством звука и размером данных для конкретной задачи.
\end{itemize}

\newpage

\thispagestyle{empty}
\begin{center}
\bfseries СОДЕРЖАНИЕ
\end{center}
\vspace{-6pt}
\tableofcontents
\clearpage

% фантомные боли для секций, чтобы не было нумераций кроме тела реферата
\section*{}
\vspace*{-2.5em}
\begin{center}\textbf{ВВЕДЕНИЕ}\end{center}
\phantomsection
\addcontentsline{toc}{section}{Введение}
Современные технологии аудиообработки играют ключевую роль в формировании цифровой медиасреды, определяя качество и доступность звукового контента. По своей сути, цифровой звук представляет собой результат дискретизации и квантования аналогового аудиосигнала, что позволяет хранить, обрабатывать и передавать его с помощью вычислительных систем. Такое преобразование лежит в основе всех современных аудиотехнологий — от стриминговых сервисов и подкастов до профессиональной звукозаписи и саунд-дизайна. Эффективная работа с аудиоконтентом требует понимания не только методов его цифрового представления, но и технологий кодирования, сжатия и последующего редактирования.

Современное состояние проблемы. Сегодня цифровой звук является неотъемлемой частью повседневной жизни, сопровождая пользователя в музыкальных сервисах, видеоконтенте, онлайн-общении и игровой индустрии. Однако за кажущейся простотой воспроизведения скрывается комплекс технологических решений. Актуальной задачей остаётся обеспечение оптимального баланса между качеством звучания, размером данных и совместимостью аудиоформатов на различных устройствах и платформах. Параллельно с развитием алгоритмов сжатия (от стандартов MP3 и AAC до современных кодеков Opus) и инструментов для редактирования (от профессиональных DAW до мобильных приложений) возрастает важность понимания психоакустических основ и принципов пространственного звука. Сфера применения расширяется от потоковой передачи музыки и создания подкастов до immersive-аудио в кино, VR-средах и интерактивных медиа, что актуализирует вопросы стандартизации, эффективного кодирования и этического использования аудиотехнологий.

Актуальность данной работы обусловлена фундаментальной ролью цифрового аудио как одного из базовых компонентов современной мультимедийной коммуникации. Умение осознанно выбирать форматы, методы обработки и средства создания аудиоконтента становится важным практическим навыком. При этом необходимо учитывать как технические ограничения (пропускная способность каналов, возможности воспроизводящих устройств), так и качественные критерии, чтобы конечный продукт соответствовал задачам и ожиданиям аудитории. Распространение иммерсивных технологий и высококачественных аудиоформатов делает изучение принципов аудиообработки необходимым для специалистов в области медиа, звукорежиссуры и цифровых коммуникаций.

Работа заключается в комплексном исследовании технологий цифрового аудио: от анализа фундаментальных принципов представления и кодирования звука до изучения практических инструментов для его обработки, ключевых форматов хранения и передачи, а также перспективных направлений развития, таких как объёмный и пространственный звук.

Цель работы: исследовать технологии создания, кодирования, обработки и воспроизведения цифрового аудиоконтента, а также проанализировать тенденции развития аудиотехнологий в контексте современных медиа и коммуникационных систем.

Задачи исследования: раскрыть принципы аналого-цифрового преобразования звукового сигнала (дискретизация, квантование) и теоретические основы данного процесса; проанализировать методы и алгоритмы сжатия аудиоданных, провести сравнительную характеристику технологий сжатия с потерями и без потерь; систематизировать знания о ключевых аудиоформатах и кодеках, определить области их эффективного применения; исследовать инструментарий и базовые техники цифрового редактирования и монтажа аудио; рассмотреть современные технологии объёмного и пространственного звука (Dolby Atmos, Ambisonics и др.) и перспективы их развития.

Методы исследования: теоретический анализ научной и технической литературы; сравнительный анализ технологий кодирования и аудиоформатов; изучение функциональных возможностей современного программного обеспечения для обработки звука; обзор актуальных стандартов и тенденций в области аудиотехнологий.

\newpage
% вручную увеличить номер секции и записать в оглавление с номером,
% чтобы в теле не печатался автоматический номер, но в tableofcontents он остался
\refstepcounter{section}
\addcontentsline{toc}{section}{\protect\numberline{\thesection} Основы представления звука в цифровых системах}
\vspace*{-2.5em}
\begin{center}\textbf{\thesection\ ОСНОВЫ ПРЕДСТАВЛЕНИЯ ЗВУКА В ЦИФРОВЫХ СИСТЕМАХ}\end{center}
\vspace*{-1.5em}
\ManualSubsection{Аналоговый звук и принцип аналого-цифрового преобразования}

Исходным объектом для цифрового представления является аналоговый звуковой сигнал, представляющий собой непрерывную механическую волну, характеризующуюся амплитудой, частотой и фазой. Фундаментальный процесс перехода от этого непрерывного представления к дискретному описывается как аналого-цифровое преобразование (АЦП). Данный процесс служит технической основой для всей современной аудиоиндустрии, начиная от телефонии и заканчивая высококачественными студийными записями. Сущность преобразования заключается в замене непрерывной функции времени последовательностью её отдельных значений, что делает возможным обработку, хранение и передачу звука цифровыми системами. Базовым методом импульсно-кодовой модуляции (ИКМ или PCM), лежащим в основе большинства цифровых аудиосистем, является стандарт G.711, разработанный Международным союзом электросвязи \cite{itu2021}.

\ManualSubsection{Дискретизация: принцип временной выборки}

Первым и ключевым этапом АЦП является дискретизация, или сэмплирование. Этот процесс заключается в измерении мгновенных значений амплитуды аналогового сигнала через равные промежутки времени. Количество таких измерений в секунду называется частотой дискретизации и измеряется в герцах (Гц). Теоретической основой корректной дискретизации без потери информации служит теорема Котельникова-Найквиста. Согласно данной теореме, для точного восстановления исходного сигнала частота дискретизации должна как минимум вдвое превышать наивысшую частоту, присутствующую в сигнале. Несоблюдение этого условия приводит к возникновению необратимых искажений, известных как алиасинг (наложение спектров). Для предотвращения алиасинга перед процессом дискретизации применяется аналоговый или цифровой фильтр нижних частот (антиалиасинговый фильтр), ограничивающий спектр входного сигнала. Типичные значения частоты дискретизации варьируются от 8 кГц для телефонной связи до 44.1 кГц для аудио-CD и 48–192 кГц для профессиональной звукозаписи. Подробное описание математических основ дискретизации представлено в руководстве по цифровой обработке сигналов \cite{smith1997dsp}.

\ManualSubsection{Квантование: измерение амплитуды}

Следующим этапом после дискретизации является квантование, в ходе которого каждое измеренное значение амплитуды округляется до ближайшего уровня из конечного набора разрешённых значений. Количество этих уровней определяется разрядностью, или глубиной квантования, которая измеряется в битах. Так, разрядность в 8 бит позволяет представить 256 возможных уровней, 16 бит — 65 536 уровней, а 24 бита — более 16 миллионов. Более высокая разрядность обеспечивает больший динамический диапазон (разница между самым тихим и самым громким звуком, который можно закодировать без искажений) и снижает уровень шума квантования — фундаментальной погрешности, возникающей из-за округления. Динамический диапазон для линейного квантования рассчитывается приблизительно как 6.02 дБ на бит, что даёт, например, около 96 дБ для 16-битного звука. Для маскировки нелинейных искажений, особенно заметных на низких уровнях сигнала, на этапе квантования может применяться техника дизеринга — добавление специального низкоуровневого шума.

\ManualSubsection{Кодирование и форматы представления}

Результатом этапов дискретизации и квантования является поток числовых значений, который затем подвергается кодированию для хранения или передачи. Наиболее прямым методом является импульсно-кодовая модуляция (PCM), где значения амплитуды просто записываются последовательно. Однако несжатый PCM-поток требует значительной полосы пропускания, например, для стандарта аудио-CD битрейт составляет 1411 кбит/с. Для оптимизации используют методы сжатия данных. Сжатие без потерь, применяемое в таких форматах, как FLAC или ALAC, позволяет уменьшить объём данных без какого-либо изменения аудиосодержания, что достигается за счёт устранения статистической избыточности \cite{xiph2024flac}. Сжатие с потерями, используемое в MP3 или AAC, основывается на психоакустических моделях и удаляет из сигнала компоненты, считаемые малозаметными для человеческого слуха, что позволяет достичь многократного сокращения размера файла.

\ManualSubsection{Практические аспекты и применение}

Выбор параметров оцифровки является компромиссом между качеством звучания, объёмом данных и вычислительными ресурсами. Для профессиональной студийной записи и последующей обработки стандартом являются параметры 24 бит / 48–96 кГц, что обеспечивает запас по динамическому диапазону и частотному разрешению. Для финального распространения музыки потребителям доминируют форматы сжатия с потерями с высоким битрейтом (AAC, Opus) на стриминговых платформах или lossless-форматы (FLAC) для аудиофильского сегмента. Современным трендом является развитие иммерсивных аудиотехнологий, таких как Dolby Atmos или Sony 360 Reality Audio, которые расширяют принципы пространственного представления звука за пределы традиционных каналов, требуя при этом новых подходов к кодированию и воспроизведению.

\newpage
\refstepcounter{section}
\addcontentsline{toc}{section}{\protect\numberline{\thesection}Методы кодирования, декодирования и передачи аудиосигналов}
\begin{center}\textbf{\thesection\ МЕТОДЫ КОДИРОВАНИЯ, ДЕКОДИРОВАНИЯ И ПЕРЕДАЧИ АУДИОСИГНАЛОВ}\end{center}
\vspace*{-1.5em}
\ManualSubsection{Принцип аналого-цифрового преобразования и импульсно-кодовая модуляция}

Ключевым процессом для представления звука в цифровых системах является аналого-цифровое преобразование (АЦП), наиболее распространённой формой которого выступает импульсно-кодовая модуляция (ИКМ или PCM). Этот процесс переводит непрерывный аналоговый сигнал в дискретную последовательность числовых значений посредством двух последовательных операций: дискретизации и квантования. Дискретизация заключается в измерении мгновенных значений амплитуды сигнала через равные интервалы времени, определяемые частотой дискретизации. Теоретической основой для выбора корректной частоты служит теорема Котельникова-Шеннона, устанавливающая, что для восстановления сигнала без потерь частота выборки должна превышать удвоенную верхнюю граничную частоту его спектра \cite{shannon1948mathematical}. Квантование присваивает каждому измеренному значению амплитуды определённый уровень из конечного набора, точность которого задаётся разрядностью. Результатом этого процесса является исходный поток PCM-данных, характеризующийся значительным объёмом, что создаёт необходимость в последующем эффективном сжатии для хранения и передачи.

\ManualSubsection{Методы сжатия аудиоданных}

Для уменьшения избыточности PCM-потока применяются алгоритмы сжатия, которые подразделяются на два фундаментальных класса. Сжатие без потерь (lossless) обеспечивает точное восстановление исходных данных после декодирования за счёт устранения статистической и перцептуальной избыточности без удаления информации. Принцип работы таких кодеков, к которым относятся FLAC или ALAC, основан на энтропийном кодировании и предсказании последующих значений сигнала на основе предыдущих. Сжатие с потерями (lossy) достигает значительно более высоких коэффициентов компрессии за счёт необратимого удаления аудиоинформации, считаемой малозаметной для человеческого восприятия. Данный подход базируется на моделях психоакустики, которые формализуют такие явления, как частотное и временное маскирование. Кодек анализирует спектр сигнала и удаляет компоненты, уровень которых находится ниже рассчитанных порогов слышимости. К этому классу относятся повсеместно распространённые форматы MP3, AAC и современный универсальный кодек Opus, спецификация которого детально описывает гибридную модель сжатия \cite{opus2024}.

\ManualSubsection{Процесс декодирования и цифро-аналогового преобразования}

Обратное преобразование цифрового аудиопотока в аналоговую форму, пригодную для воспроизведения акустическими системами, включает этап декодирования (декомпрессии) и цифро-аналогового преобразования (ЦАП). Для сжатых данных сначала выполняется декомпрессия. В случае lossless-кодеков результатом является исходный PCM-сигнал, идентичный закодированному. При декодировании lossy-форматов происходит не восстановление, а синтез PCM-сигнала на основе сохранённых параметров, что приводит к получению аппроксимации исходного звука. Последующий процесс ЦАП заключается в преобразовании последовательности дискретных числовых значений в ступенчатый аналоговый сигнал. Ключевой задачей на этом этапе является применение сглаживающего (восстанавливающего) фильтра нижних частот, который подавляет спектральные артефакты, возникающие из-за дискретной природы сигнала, и формирует непрерывную выходную волну.

\ManualSubsection{Технологии передачи цифрового аудио}

После кодирования и сжатия цифровой аудиопоток передаётся по различным каналам связи. В проводных сетях, включая интернет, для потоковой передачи часто используются адаптивные протоколы (HLS, DASH), которые динамически подстраивают битрейт аудио под доступную пропускную способность сети, используя эффективные кодеки, такие как AAC или Opus. Для передачи несжатого или сжатого многоканального звука в бытовой технике применяются специализированные интерфейсы: S/PDIF (оптический или коаксиальный) и HDMI. В области беспроводной передачи доминирующей технологией является Bluetooth, в рамках которой качество звука определяется используемым аудиокодеком. Базовым и обязательным кодеком в стеке Bluetooth является SBC (Subband Codec), в то время как для повышения качества и снижения задержки используются проприетарные решения, такие как aptX, aptX HD или LDAC, спецификации которых являются частью общего стандарта Bluetooth \cite{bluetooth2024}. Альтернативой выступает передача по Wi-Fi, которая благодаря большей пропускной способности позволяет транслировать аудио высокого разрешения без дополнительного сжатия с потерями.

\newpage
\refstepcounter{section}
\addcontentsline{toc}{section}{\protect\numberline{\thesection}Отличие сжатия без потерь и с потерями при хранении аудио}
\begin{center}\textbf{\thesection\ ОТЛИЧИЕ СЖАТИЯ БЕЗ ПОТЕРЬ И С ПОТЕРЯМИ ПРИ ХРАНЕНИИ АУДИО}\end{center}
\vspace*{-1.5em}
\ManualSubsection{Мотивация и принципы сжатия аудиоданных}

Основной задачей сжатия аудиоданных является оптимизация объёма файлов для хранения и передачи без критического снижения воспринимаемого качества. Исходный несжатый цифровой аудиопоток, например в формате PCM, характеризуется значительной избыточностью, что делает прямое его использование в потребительских приложениях неэффективным. Два фундаментально различающихся подхода к решению этой проблемы — сжатие без потерь (lossless) и сжатие с потерями (lossy) — сформировали современную экосистему аудиоформатов, предлагая различные компромиссы между точностью воспроизведения и степенью уменьшения размера данных.

\ManualSubsection{Сжатие без потерь}

Принцип сжатия без потерь заключается в алгоритмическом устранении статистической избыточности из аудиопотока с гарантией последующего точного восстановления исходных данных. Данный процесс можно сравнить с архивацией: кодер анализирует последовательность сэмплов, предсказывает последующие значения и кодирует разницу между предсказанием и реальными данными методами энтропийного кодирования, например, по алгоритму Хаффмана. Типичные коэффициенты сжатия для аудио составляют от 1.5:1 до 3:1. Наиболее распространёнными форматами данного класса являются FLAC (Free Lossless Audio Codec) и ALAC (Apple Lossless Audio Codec), последний из которых, будучи открытым кодеком от Apple, часто используется в экосистеме данной компании \cite{alac_apple}. Основным преимуществом lossless-подхода является полное сохранение аудиоинформации, что делает его незаменимым для архивного хранения, профессиональной звукозаписи и использования аудиофильским сообществом.

\ManualSubsection{Сжатие с потерями}

В отличие от lossless-подхода, сжатие с потерями обеспечивает радикально большее уменьшение размера файла (до 10:1 и выше) за счёт необратимого удаления части аудиоинформации, считаемой малозаметной для человеческого слуха. Этот процесс основывается на психоакустических моделях, которые формализуют явления частотного и временного маскирования. Кодер проводит частотный анализ сигнала и удаляет спектральные компоненты, уровень которых находится ниже расчётных порогов слышимости. Исторически основополагающим стандартом стал MPEG-1 Audio Layer III (MP3) \cite{iso11172_3}. Более совершенным и современным форматом является Advanced Audio Coding (AAC), который при сопоставимом битрейте обеспечивает лучшее качество звучания за счёт улучшенных алгоритмов \cite{iso14496_3}. Также следует отметить открытый формат Ogg Vorbis, использующий аналогичные принципы и представленный в соответствующей спецификации \cite{vorbis_spec}.

\ManualSubsection{Сравнительный анализ и области применения}

Ключевое отличие между двумя подходами заключается в фундаментальном компромиссе между акустической точностью и эффективностью сжатия. Lossless-кодеки гарантируют битовую идентичность декодированного сигнала оригиналу, что определяет их применение в сценариях, где неприемлема любая деградация качества. В свою очередь, lossy-форматы достигают высокой степени сжатия, достаточной для потоковой передачи и хранения обширных музыкальных коллекций на устройствах с ограниченной памятью, при этом качество звучания на высоких битрейтах (256–320 кбит/с) для большинства слушателей является субъективно близким к оригиналу. Выбор конкретного формата и параметров сжатия, таких как постоянный (CBR) или переменный (VBR) битрейт, диктуется требованиями к качеству, доступными ресурсами хранения и пропускной способностью каналов связи. В современных медиасервисах, таких как Spotify, часто используется формат Ogg Vorbis, в то время как Apple Music предлагает опциональное прослушивание в lossless-качестве наряду со стандартным lossy-потоком на основе AAC.

\newpage
\refstepcounter{section}
\addcontentsline{toc}{section}{\protect\numberline{\thesection}Основные форматы звуковых файлов и области их применения}
\begin{center}\textbf{\thesection\ ОСНОВНЫЕ ФОРМАТЫ ЗВУКОВЫХ ФАЙЛОВ И ОБЛАСТИ ИХ ПРИМЕНЕНИЯ}\end{center}
\vspace*{-1.5em}
\ManualSubsection{Классификация и принципы работы аудиоформатов}

Современные форматы цифрового аудио подразделяются на два фундаментальных класса в зависимости от алгоритма хранения данных: без потерь (lossless) и с потерями (lossy). Форматы без потерь, такие как FLAC (Free Lossless Audio Codec), ALAC (Apple Lossless Audio Codec) или несжатый WAV/AIFF, обеспечивают точное, байт-в-байт соответствие исходному PCM-потоку. Их алгоритмы, подобные применяемым в FLAC, устраняют статистическую избыточность аудиоданных, используя предсказание и энтропийное кодирование, что позволяет уменьшить объём файла примерно на 30–50\%. Форматы с потерями, наиболее известными из которых являются MP3 и AAC (Advanced Audio Coding), достигают радикально большего сжатия (до 90\%) за счёт необратимого удаления аудиоинформации, считаемой малозаметной для слуха на основе психоакустических моделей. Эти модели формализуют явления частотного и временного маскирования, позволяя кодеку отбрасывать спектральные компоненты, заглушаемые более громкими звуками в тот же момент времени или на близких частотах. Технически данный процесс реализуется через преобразование сигнала в частотную область (например, с помощью MDCT — Modified Discrete Cosine Transform), анализ с помощью психоакустической модели, квантование и последующее энтропийное кодирование.

\ManualSubsection{Сравнительный анализ форматов и влияние битрейта}

Ключевым параметром, определяющим компромисс между качеством и размером файла в lossy-форматах, является битрейт — скорость потока данных, измеряемая в килобитах в секунду (кбит/с). Для формата MP3 битрейт в 128 кбит/с считается минимально приемлемым для музыки, 192 кбит/с обеспечивает хорошее качество для большинства слушателей, а 320 кбит/с приближается к субъективно «прозрачному» уровню, когда отличия от оригинала становятся малозаметными на бытовом оборудовании. Более современный формат AAC при сопоставимом битрейте обычно обеспечивает лучшее качество звучания благодаря усовершенствованным алгоритмам преобразования и психоакустической модели. Важным аспектом является и структура файла, включающая не только аудиоданные, но и метаданные (теги). Форматы используют разные системы для их хранения: MP3 опирается на теги ID3, FLAC — на Vorbis comments, а AAC в контейнере MP4/M4A — на структуру атомов MP4. Поддержка этих систем метаданных варяируется в разных устройствах и программных плеерах.

\ManualSubsection{Области применения и выбор формата}

Выбор оптимального аудиоформата диктуется конкретным сценарием использования, балансируя между требованиями к качеству, объёмом данных и совместимостью. В профессиональной музыкальной production, на этапах записи, сведения и мастеринга, используются исключительно несжатые (WAV, AIFF) или lossless-форматы (FLAC, ALAC) с высокими параметрами (24/32 бит, 48–192 кГц). Это обеспечивает максимальное качество и «запас» для неразрушающей обработки в цифровых аудио рабочих станциях (DAW), таких как Pro Tools, Cubase или Ableton Live \cite{daw_wikipedia}. Для конечного распространения музыки, особенно через интернет-стриминг, доминируют lossy-форматы. Сервисы вроде Spotify (Ogg Vorbis) или Apple Music (AAC, с опцией lossless) выбирают их из-за критически важной экономии пропускной способности и места на серверах при сохранении качества, достаточного для большинства пользователей. Для архивного долгосрочного хранения мастер-записей рекомендуется использовать открытые lossless-форматы (например, FLAC) или стандартизированные несжатые форматы (Broadcast WAV), сопровождая их контрольными суммами и соблюдая правило резервного копирования 3-2-1.

\ManualSubsection{Инструменты обработки и преобразования форматов}

Работа с разнообразными аудиоформатами требует специализированного программного инструментария. Универсальным решением для конвертации, базового редактирования и анализа медиафайлов является пакет FFmpeg, документация которого охватывает поддержку огромного количества кодеков и контейнеров \cite{ffmpeg_doc}. Для неразрушающего редактирования аудио, чистки записей, наложения эффектов и работы с многодорожечными проектами используются цифровые аудио рабочие станции (DAW) и аудиоредакторы. Современные DAW представляют собой комплексные программные среды, централизованно управляющие записью, редактированием, виртуальными инструментами и сведением аудиотреков \cite{daw_wikipedia}. В свою очередь, для решения более узких задач, таких как шумоподавление, анализ спектра или простой монтаж, подходят специализированные аудиоредакторы. Таким образом, эффективная работа в современной аудиосфере подразумевает понимание не только характеристик форматов, но и возможностей инструментария для их практического использования.

\newpage
\refstepcounter{section}
\addcontentsline{toc}{section}{\protect\numberline{\thesection}Основы редактирования и монтажа аудио в цифровой среде}
\begin{center}\textbf{\thesection\ ОСНОВЫ РЕДАКТИРОВАНИЯ И МОНТАЖА АУДИО В ЦИФРОВОЙ СРЕДЕ}\end{center}
\vspace*{-1.5em}
\ManualSubsection{Фундаментальные принципы цифрового аудио и программной среды}

Процесс редактирования и монтажа в цифровой аудиосреде базируется на представлении звука в виде дискретной последовательности сэмплов. Такое представление является результатом аналого-цифрового преобразования (АЦП), ключевыми параметрами которого выступают частота дискретизации и разрядность (битовая глубина). Частота дискретизации определяет верхнюю граничную частоту воспроизводимого спектра (согласно теореме Котельникова-Найквиста), а разрядность задаёт динамический диапазон и уровень шума квантования. Исходные данные для профессионального редактирования, как правило, хранятся в несжатых форматах (WAV, AIFF) или в форматах сжатия без потерь (FLAC, ALAC), что исключает деградацию качества при многократных операциях. Основным программным инструментом для работы выступает цифровая аудио рабочая станция (DAW). Современные DAW, будь то профессиональные решения (Avid Pro Tools, Steinberg Cubase, Apple Logic Pro) или бесплатные пакеты (Audacity, GarageBand), обеспечивают комплексную среду, включающую монтажную область для расположения клипов на дорожках, микшер для регулировки уровней и панорамирования, а также систему для подключения виртуальных инструментов и плагинов эффектов.

\ManualSubsection{Базовые техники неразрушающего редактирования и обработки}

Технологический процесс цифрового аудиомонтажа строится на принципе неразрушающего (non-destructive) редактирования, при котором все операции применяются к данным в реальном времени, не изменяя исходные аудиофайлы. Базовый цикл работы включает импорт и организацию медиаматериалов по дорожкам, соответствующим различным компонентам проекта (дикторский голос, музыка, шумы). Фундаментальные монтажные операции выполняются с помощью инструментов «выделение» и «разрезание», что позволяет точно позиционировать, обрезать и удалять фрагменты. Для создания плавных переходов между аудиоклипами применяется функция склейки с автоматическим наложением короткого перекрёстного затухания (crossfade). Параллельно с монтажом осуществляется первичная обработка звука: регулировка относительного уровня громкости (gain staging), балансировка стереопанорамы (panning) и частотная коррекция с помощью эквалайзера для устранения проблемных частот или выделения значимых элементов микса. Важной частью процесса является также очистка записи от артефактов с использованием де-эссеров для подавления свистящих согласных, декликеров для удаления щелчков и алгоритмов шумоподавления для селективного ослабления постоянного фона.

\ManualSubsection{Творческая обработка с использованием эффектов и пространственная организация}

Качественный аудиомонтаж выходит за рамки базовой компоновки и включает применение творческих эффектов для формирования определённого характера и пространства звучания. Ключевыми инструментами здесь являются динамическая обработка и эффекты, основанные на времени задержки сигнала. Динамический процессор (компрессор) используется для выравнивания громкостного диапазона, что делает тихие элементы более слышимыми, а пики — контролируемыми. Для добавления ощущения пространства и объёма применяется реверберация, имитирующая акустику конкретного помещения, и задержка (delay), создающая эффект эха. В контексте современных медиапродуктов, особенно в кино, видеоиграх и иммерсивных музыкальных релизах, возрастает важность работы с многоканальными и объектно-ориентированными форматами пространственного звука, такими как Dolby Atmos \cite{dolby_atmos_guide} и DTS:X \cite{dts_x_standard}. Эти технологии позволяют позиционировать звуковые объекты в трёхмерном пространстве вокруг слушателя, что требует от монтажёра и звукорежиссёра новых навыков работы в специализированных DAW, поддерживающих соответствующие форматы рендеринга, как это реализовано в экосистеме Apple \cite{apple_spatial_audio}.

\ManualSubsection{Финальная стадия: сведение и мастеринг}

Завершающей стадией производства является мастеринг — процесс финальной оптимизации и подготовки сведённой стереофайлы (или многоканального микса) к целевому формату распространения. Цель мастеринга — обеспечить конкурентную громкость, частотный баланс и общее качество звучания, соответствующее отраслевым стандартам, а также устранить возможные технические дефекты, незамеченные на этапе сведения. Для этого используются специализированные высокоточные инструменты. Мультибандовый компрессор позволяет независимо контролировать динамику в различных частотных диапазонах. Линейный фазовый эквалайзер осуществляет тонкую коррекцию тембра без внесения фазовых искажений. Стерео процессор (имиджер) управляет шириной звуковой сцены. Финальное повышение интегральной громкости и ограничение пиков для предотвращения цифрового клиппинга выполняется с помощью лимитера или максимайзера. Результатом профессионального мастеринга является мастер-файл, готовый для последующего тиражирования на физические носители (CD, винил) или кодирования в форматы сжатия (MP3, AAC) для цифровых платформ и стриминговых сервисов.

\newpage
\refstepcounter{section}
\addcontentsline{toc}{section}{\protect\numberline{\thesection}Перспективы развития технологий объемного и пространственного звука}
\begin{center}\textbf{\thesection\ ПЕРСПЕКТИВЫ РАЗВИТИЯ ТЕХНОЛОГИЙ ОБЪЕМНОГО И ПРОСТРАНСТВЕННОГО ЗВУКА}\end{center}
\vspace*{-1.5em}
\ManualSubsection{Эволюция от канального к объектно-ориентированному аудио}

Современные технологии объемного звука представляют собой эволюционный скачок от традиционных канальных систем, к объектно-ориентированным и сценарным форматам. Ключевое концептуальное отличие заключается в том, что вместо фиксированного набора аудиодорожек, привязанных к конкретным колонкам, звук описывается как набор независимых объектов, наделенных метаданными о позиции, траектории движения и других атрибутах в трехмерном пространстве. Такая архитектура, лежащая в основе стандартов Dolby Atmos и MPEG-H 3D Audio \cite{mpeg_immersive_audio}, обеспечивает гибкость воспроизведения: финальный рендеринг звуковой сцены происходит в реальном времени на стороне декодера, адаптируясь под конкретную акустическую конфигурацию — от многоканального кинотеатрального комплекса до системы домашнего кинотеатра, звуковой панели или даже стереонаушников. Это преодолевает главное ограничение канального звука — жесткую привязку к предопределённой геометрии системы.

\ManualSubsection{Ключевые технологии создания иммерсивного опыта}

Реализация убедительного трехмерного аудиоопыта базируется на нескольких взаимодополняющих технологических подходах. Для воспроизведения в наушниках ключевую роль играет биноуральный рендеринг, который с высокой точностью моделирует физические процессы распространения звука от источника до барабанных перепонок слушателя, учитывая индивидуальные особенности строения ушной раковины (HRTF — Head-Related Transfer Function). Совмещение этого метода с отслеживанием положения головы в пространстве (head tracking) позволяет «зафиксировать» виртуальные звуковые источники в окружающей среде, значительно усиливая иллюзию реальности и стабильность звуковой сцены. Для громкоговорящих систем актуальны технологии адаптивной акустики и лучеформирования (beamforming). Смарт-колонки и саундбары, оснащенные массивами динамиков, с помощью цифровой обработки сигналов формируют направленные звуковые лучи, которые отражаются от стен помещения для создания виртуальных источников звука с эффектом высоты, что устраняет необходимость в физической установке потолочных динамиков. Данные подходы подробно рассматриваются в исследованиях по адаптивным акустическим системам \cite{ieee_adaptive_audio}.

\ManualSubsection{Применение в медиаиндустрии и будущие направления}

Внедрение пространственного звука активно трансформирует ключевые сегменты медиаиндустрии. В кинематографе и производстве музыкального контента объектно-ориентированный звук предоставляет авторам беспрецедентный творческий контроль над звуковой перспективой и движением. В индустрии видеоигр и виртуальной реальности (VR) точная 3D-аудиолокализация является не просто элементом погружения, а критически важным игровым механизмом, влияющим на геймплей. Актуальным трендом является развитие стандартов для потоковой передачи иммерсивного аудио, что отражено в дорожных картах профессиональных сообществ \cite{aes_spatial_audio_roadmap}. Перспективным направлением считается персонализация пространственного звука, где системы будут адаптировать аудиопараметры под индивидуальные анатомические особенности слуха пользователя и акустику конкретного помещения. Наиболее дальновидные сценарии связывают эволюцию пространственного аудио с развитием концепций метавселенных и смешанной реальности (XR), где звук будет выступать фундаментальным слоем для создания убедительных, устойчивых и социальных виртуальных сред, что отмечается в исследованиях на эту тему \cite{meta_audio_arxiv}. Таким образом, пространственный звук эволюционирует от специализированной технологии для домашних кинотеатров в базовую, повсеместно ожидаемую составляющую цифрового медиаопыта.

\newpage
\section*{}
\vspace*{-2.5em}
\begin{center}\textbf{ЗАКЛЮЧЕНИЕ}\end{center}
\phantomsection
\addcontentsline{toc}{section}{Заключение}

Цифровые аудиотехнологии, основанные на фундаментальных процессах дискретизации, квантования и кодирования, сформировали техническую инфраструктуру всей современной звуковой индустрии. В ходе проведённого исследования были систематизированы базовые принципы представления звука в цифровых системах, проанализированы методы кодирования и передачи аудиосигналов, а также детально рассмотрена классификация аудиоформатов на основе алгоритмов сжатия — без потерь (lossless) и с потерями (lossy). Каждый из этих аспектов представляет собой результат компромисса между такими ключевыми параметрами, как качество звучания, размер данных, вычислительная сложность и совместимость с воспроизводящим оборудованием. Понимание этой взаимосвязи позволяет осознанно выбирать инструменты и форматы для решения конкретных задач — от студийной звукозаписи до массового потокового вещания.

Особое внимание в работе уделено практическим аспектам работы со звуком в цифровой среде, включая основы редактирования и монтажа в специализированных рабочих станциях (DAW). Эти инструменты, предоставляющие возможности неразрушающей обработки, применения эффектов и точного контроля над всеми параметрами звука, являются ключевым звеном в цепочке создания современного аудиоконтента. Итоговым, завершающим этапом этого процесса выступает мастеринг, цель которого — финальная оптимизация и унификация звучания для достижения коммерческого качества и соответствия отраслевым стандартам. Таким образом, целостный технологический цикл — от оцифровки исходного сигнала до финального мастеринга — обеспечивает создание профессионального звукового продукта.

Наиболее динамично развивающимся направлением, определяющим будущее аудиоиндустрии, являются технологии объёмного и пространственного звука (Immersive Audio). Переход от канальных систем к объектно-ориентированным форматам, таким как Dolby Atmos и MPEG-H, а также развитие биноурального рендеринга и адаптивной акустики, знаменуют смену парадигмы — от простого воспроизведения звука к моделированию целостной трёхмерной акустической среды. Эти технологии выходят за рамки сферы развлечений, находя применение в телекоммуникациях, образовании и индустрии виртуальной реальности, где точная пространственная аудиолокализация становится критически важной для полного погружения и эффективного взаимодействия.

Таким образом, современные аудиотехнологии представляют собой сложную, многоуровневую экосистему, объединяющую фундаментальные принципы обработки сигналов, передовые алгоритмы сжатия данных, мощный инструментарий для творческой работы и перспективные стандарты иммерсивного звука. Для специалиста в области медиа и звукорежиссуры владение этими знаниями является неотъемлемой компетенцией, позволяющей не только эффективно работать с существующими инструментами, но и критически оценивать тенденции развития, адаптируясь к стремительным изменениям в технологическом ландшафте цифрового аудио.

\newpage
\begin{center}\textbf{СПИСОК СОКРАЩЕНИЙ}\end{center}
\phantomsection
\addcontentsline{toc}{section}{Список сокращений}

\vspace{1em}

{
\small
\noindent
\begin{tabular}{|l|p{0.75\textwidth}|}
\hline
\textbf{Сокращение} & \multicolumn{1}{c|}{\textbf{Расшифровка}} \\
\noalign{\hrule height 1.2pt}
\hline
АЦП & Аналого-цифровой преобразователь (устройство или процесс преобразования непрерывного аналогового сигнала в цифровую форму)\\
\hline
ИКМ & Импульсно-кодовая модуляция (Pulse-Code Modulation, PCM — метод цифрового представления аналогового сигнала, основанный на дискретизации и квантовании)\\
\hline
PCM & Pulse-Code Modulation (Импульсно-кодовая модуляция; базовый формат несжатого цифрового аудио)\\
\hline
AAC & Advanced Audio Coding (Усовершенствованное кодирование аудио; современный формат сжатия с потерями, более эффективный, чем MP3)\\
\hline
FLAC & Free Lossless Audio Codec (свободный аудиокодек без потерь)\\
\hline
ALAC & Apple Lossless Audio Codec (Аудиокодек Apple без потерь; проприетарный формат сжатия без потерь в экосистеме Apple)\\
\hline
CBR & Constant Bit Rate (Постоянный битрейт; метод кодирования, при котором битрейт не меняется на протяжении файла)\\
\hline
DAW & Digital Audio Workstation (Цифровая аудио рабочая станция; программное обеспечение для записи, редактирования и производства аудио)\\
\hline
HTML & HyperText Markup Language (язык гипертекстовой разметки)\\
\hline
HRTF & Head-Related Transfer Function (Функция передачи, связанная с головой; математическая модель, описывающая, как уши человека фильтруют звук из определённой точки в пространстве)\\
\hline
MDCT & Modified Discrete Cosine Transform (Модифицированное дискретное косинусное преобразование; математическое преобразование, широко используемое в аудиокодеках для перевода сигнала в частотную область)\\
\hline
UX & User Experience (пользовательский опыт)\\
\hline
MP3 & MPEG-1/MPEG-2 Audio Layer III (аудиоформат с потерями)\\
\hline
MP4 & MPEG-4 Part 14 (цифровой мультимедийный формат-контейнер)\\
\hline
VBR & Variable Bit Rate (Переменный битрейт; метод кодирования, при котором битрейт динамически меняется в зависимости от сложности аудиоданных)\\
\hline
\end{tabular}
}

\clearpage
\nocite{*}
\printbibliography[heading=bibliography]

\end{document}
